\appendix
% \section{Drafted ROS2 System Architecture}
% \label{sec:AppA}

% Figure~\ref{fig:ros2_arch} shows the drafted ROS~2 node architecture for the full system. Red dashed nodes and topics are planned additions for Term~2. Table~\ref{tab:topics} shows the ROS~2 topics and their descriptions.

% \begin{figure}[H]
% \centering
% \usetikzlibrary{shapes.geometric, arrows.meta, positioning, backgrounds, fit}
% \resizebox{\textwidth}{!}{%
% \begin{tikzpicture}[
%     ros/.style={rectangle, rounded corners=3pt, draw=black, thick,
%                 fill=blue!10, minimum width=3.0cm, minimum height=1.0cm,
%                 font=\small, text centered},
%     sim/.style={rectangle, rounded corners=3pt, draw=gray!60, thick,
%                 fill=gray!10, minimum width=3.0cm, minimum height=1.0cm,
%                 font=\small, text centered},
%     planned/.style={rectangle, rounded corners=3pt, draw=red!60, thick,
%                 fill=red!5, minimum width=3.0cm, minimum height=1.0cm,
%                 font=\small, text centered, dashed},
%     tl/.style={font=\footnotesize, fill=white, inner sep=1pt},
%     arr/.style={-Stealth, thick},
%     darr/.style={-Stealth, thick, dashed, red!60},
% ]

% %% ROW 1 — Central controller
% \node[ros] (cc) {\textbf{central\_controller}};
% \node[planned, right=4.5cm of cc] (ta) {tension\_allocator$^*$};

% %% ROW 2 — Per-drone controllers
% \node[ros, below left=1.4cm and 4.5cm of cc]  (c0) {controller\_0};
% \node[ros, below left=1.4cm and 0.8cm of cc]  (c1) {controller\_1};
% \node[ros, below right=1.4cm and 0.8cm of cc] (c2) {controller\_2};
% \node[ros, below right=1.4cm and 4.5cm of cc] (c3) {controller\_3};

% %% ROW 3 — Mocap nodes
% \node[ros, below=1.2cm of c0] (m0) {mocap\_0};
% \node[ros, below=1.2cm of c1] (m1) {mocap\_1};
% \node[ros, below=1.2cm of c2] (m2) {mocap\_2};
% \node[ros, below=1.2cm of c3] (m3) {mocap\_3};

% %% Payload listener (centred)
% \node[ros, below=1.4cm of cc] (psl) {payload\_state\_listener};
% \node[planned, below=0.8cm of psl] (ps) {/payload/state$^*$};

% %% fleet/step: central to controllers
% \draw[<->] (cc) -- node[tl, midway]{\texttt{/fleet/step}} (c0);
% \draw[<->] (cc) -- (c1);
% \draw[<->] (cc) -- (c2);
% \draw[<->] (cc) -- (c3);

% %% planned tension allocator
% \draw[darr] (ta) -- node[tl]{\texttt{/fleet/tension\_refs}$^*$} (cc);

% %% command: mocap to controllers
% \draw[arr] (m0) -- node[tl, right]{\texttt{/drone\_N/command}} (c0);
% \draw[arr] (m1) -- (c1);
% \draw[arr] (m2) -- (c2);
% \draw[arr] (m3) -- (c3);

% %% payload state to central
% \draw[arr] (psl) -- node[tl, right]{\texttt{/payload\_state\_publisher}} (cc);
% \draw[darr] (ps) -- (psl);

% %% Ground computer bounding box
% \begin{scope}[on background layer]
%   \node[draw=blue!40, fill=blue!3, rounded corners=8pt, inner sep=14pt,
%         fit=(cc)(ta)(c0)(c3)(m0)(m3)(psl)(ps),
%         label={[font=\small\bfseries, blue!70]above:Ground Computer (Offboard)}] (gcbox) {};
% \end{scope}

% %% -------------------------------------------------------
% %% SIMULATION LAYER
% %% -------------------------------------------------------
% \node[sim, below=1.6cm of m0] (pb0) {pose\_bridge\_0};
% \node[sim, below=1.6cm of m1] (pb1) {pose\_bridge\_1};
% \node[sim, below=1.6cm of m2] (pb2) {pose\_bridge\_2};
% \node[sim, below=1.6cm of m3] (pb3) {pose\_bridge\_3};

% \node[sim, below=1.0cm of pb0] (gz0) {Gazebo drone\_0};
% \node[sim, below=1.0cm of pb1] (gz1) {Gazebo drone\_1};
% \node[sim, below=1.0cm of pb2] (gz2) {Gazebo drone\_2};
% \node[sim, below=1.0cm of pb3] (gz3) {Gazebo drone\_3};

% %% pose bridge chain
% \draw[arr] (gz0) -- (pb0);
% \draw[arr] (gz1) -- node[tl, right]{\texttt{/model/x3\_droneN/pose}} (pb1);
% \draw[arr] (gz2) -- (pb2);
% \draw[arr] (gz3) -- (pb3);

% \draw[arr] (pb0) -- (m0);
% \draw[arr] (pb1) -- (m1);
% \draw[arr] (pb2) -- (m2);
% \draw[arr] (pb3) -- (m3);

% %% Simulation bounding box
% \begin{scope}[on background layer]
%   \node[draw=gray!40, fill=gray!5, rounded corners=8pt, inner sep=14pt,
%         fit=(pb0)(pb3)(gz0)(gz3),
%         label={[font=\small\bfseries, black!60]below:Simulation Layer (replaced by the motion capture system and hardware in experiments)}] {};
% \end{scope}

% %% Dividing line
% \draw[thick, dashed, black!50]
%     ([yshift=-0.8cm]m0.south) -- ([yshift=-0.8cm]m3.south);

% \end{tikzpicture}
% }
% \caption{ROS~2 node architecture.}
% \label{fig:ros2_arch}
% \end{figure}

% \begin{table}[H]
% \centering
% \caption{Key ROS~2 topics in the multi-drone payload system. $^*$Planned for Term~2.}
% \label{tab:topics}
% \small
% \begin{tabular}{llp{4.0cm}}
% \hline
% \textbf{Topic} & \textbf{Description} \\
% \hline
% \texttt{/fleet/step}  & Synchronised trajectory step index at 30\,Hz \\
% \texttt{/drone\_N/command} & Desired attitude and thrust setpoint for drone $N$ \\
% \texttt{/drone\_N/motion\_capture\_state} & Pose and velocity estimate at 120\,Hz \\
% \texttt{/drone\_N/ELRSCommand} &  ELRS roll/pitch/throttle/yaw commands \\
% \texttt{/drone\_N/arming\_state\_feedback} &  Arming status from flight controller \\
% \texttt{/drone\_N/telemetry} &  Flight controller telemetry stream \\
% \texttt{/model/x3\_droneN/pose} &  Ground-truth drone pose from simulation \\
% \texttt{/model/payload/pose}  & Ground-truth payload pose from simulation \\
% \texttt{/x3\_droneN/gazebo/command/motor\_speed} & Motor speed commands to Gazebo actuators \\
% \texttt{/mpc\_plan} &  MPC trajectory plan for visualisation \\
% \texttt{/actual\_path} &  Recorded actual flight path for visualisation \\
% \hline
% \texttt{/payload/state}$^*$ & Full payload pose and velocity for tension allocator \\
% \texttt{/fleet/tension\_refs}$^*$ &  Per-drone tension reference vectors \\
% \texttt{/fleet/agent\_status}$^*$ &  Active agent reference for join/leave management \\
% \texttt{/docking/attachment\_event}$^*$  & Attachment confirmation triggering join transition \\
% \hline
% \end{tabular}
% \end{table}

% \section{Risk Assessment}
% \label{sec:AppB}

% \begin{table}[H]
% \centering
% \caption{Project Risk Assessment}
% \label{tab:risk_assessment}
% \renewcommand{\arraystretch}{1.4}
% \begin{tabularx}{\textwidth}{p{4.2cm} c c p{6.8cm}}
% \hline
% \textbf{Risk} & \textbf{Likelihood} & \textbf{Impact} & \textbf{Mitigation} \\
% \hline
% Docking controller or collaborating student integration unavailable or delayed & Medium & Medium & Proceed with agent removal only for hardware validation with agent addition demonstrated in simulation. Result remains independently publishable. \\
% \hline
% Motion capture latency or calibration errors during hardware experiments & Medium & Medium & Motion capture emulator validated at 120\,Hz in simulation before hardware phase. Early exposure to hardware to understand the limitations. \\
% \hline
% Gazebo real-time factor limits simulation iteration speed & High & Low & Reduce the complexity of the world and run the controller at the simulation time. \\
% \hline
% Hardware unavailable during T3 experiment windows & Low & High & Complete all simulation validation before hardware phase. Results can be taken from simulated experiments. \\
% \hline
% Drone damage or failure & Low & Medium & Extensive testing in simulation before hardware flights. \\
% \hline
% \end{tabularx}
% \end{table}
% ---------------------------------------------------------------------------------------------
% Appendices drafted 4 Oct 2026 (code repository state: branch week4, packages as renamed 3 Oct).
% ---------------------------------------------------------------------------------------------

\section{Software architecture}
\label{sec:app_software}

Table~\ref{tab:app_nodes} lists the nodes of a flight and Table~\ref{tab:app_topics} the main
topics between them. The same nodes run in simulation and on the rig; only the input and output
nodes differ (the Gazebo bridges and emulators against the radio and motion-capture nodes).

\begin{table}[H]
\centering
\small
\caption{Nodes of a flight.}
\label{tab:app_nodes}
\begin{tabularx}{\textwidth}{llX}
\toprule
Package & Node & Role \\
\midrule
\texttt{mpc\_planner} & \texttt{mpc\_planner} & Load--cable planner (10\,Hz): rod sweep, pretension, lift, trajectory, LAND; one reference per UAV \\
\texttt{dissipative\_planner} & \texttt{dissipative\_planner} & The planner plus reconfiguration: OCP resize for detach and attach, the dissipative network, the approach profile \\
\texttt{tracker} & \texttt{tracker\_}$i$ & Per-UAV tracking MPC (50\,Hz), or the velocity loop \\
\texttt{fleet\_manager} & \texttt{fleet\_manager} & ARM, TAKEOFF, LAND, DISARM, ESTOP; fleet-wide abort; the pre-takeoff gate \\
\texttt{drone\_communication} & \texttt{elrs\_interface} & One per UAV: CRSF frames to the radio at 100\,Hz, telemetry back, arming confirmed by the flight controller, magnet channel \\
\texttt{drone\_communication} & \texttt{motion\_capture\_publisher} & OptiTrack rigid bodies to per-body pose and velocity topics \\
\texttt{drone\_magnet} & magnet manager, \texttt{elrs\_mux} & Weld detection and the authority hand-over at an attach \\
\texttt{drone\_visualisation} & panel, \texttt{fleet\_viz} & Operator panel (ARM, TAKEOFF, LAND, DISARM, magnets) and RViz scene \\
\texttt{simulation\_communication} & bridges, emulators & Gazebo only: flight controller and motion-capture emulators, clock \\
\bottomrule
\end{tabularx}
\end{table}

\begin{table}[H]
\centering
\small
\caption{Main topics.}
\label{tab:app_topics}
\begin{tabularx}{\textwidth}{lX}
\toprule
Topic & Content \\
\midrule
\texttt{/drone\_}$i$\texttt{/motion\_capture\_state}, \texttt{/payload/...} & Pose and velocity of each UAV and of the payload \\
\texttt{/drone\_}$i$\texttt{/reference\_trajectory} & The planner's horizon for UAV $i$: node count, time step, then per node position, velocity, acceleration feedforward and planned cable acceleration (twelve numbers per node) \\
\texttt{/drone\_}$i$\texttt{/ELRSCommand} & The tracker's command: armed flag, throttle, three rate sticks, auxiliary channels \\
\texttt{/drone\_}$i$\texttt{/telemetry} & Battery voltage, flight mode, link quality \\
\texttt{/drone\_}$i$\texttt{/fc\_arm\_state} & The flight controller's own arm state and the reason for a refusal \\
\texttt{/drone\_}$i$\texttt{/magnet} & Magnet on or off \\
\texttt{/fleet/command} & ARM, TAKEOFF, LAND, DISARM, ESTOP \\
\texttt{/fleet/detach}, \texttt{/magnet/command} & DETACH of UAV $i$; ATTACH (magnet on) \\
\texttt{/fleet/abort} & Fleet-wide disarm on any fault \\
\texttt{/fleet/landed}, \texttt{/fleet/cable\_ff\_active} & Planner to fleet: landed; cable feedforward on \\
\bottomrule
\end{tabularx}
\end{table}

\section{Simulation and rig parameters}
\label{sec:app_params}

\begin{table}[H]
\centering
\small
\caption{Parameters of the rig twin and the rig.}
\label{tab:app_params}
\begin{tabularx}{\textwidth}{>{\raggedright\arraybackslash}Xcc}
\toprule
Parameter & Rig twin & Rig \\
\midrule
UAV mass (kg) & 0.55 & 0.55 \\
rod length (m) & 0.55 & 0.55 \\
attachment radius (m) & 0.225 & 0.225 \\
rod pivot below the UAV centre (m) & 0.04 & 0.04 \\
thrust map & affine & affine \\
thrust gain $k_T$ & 35.2 & 35.2 \\
per-UAV thrust trim & off & off \\
throttle cap & 0.8 & 0.8 \\
lift rate (m/s) & 0.1 & 0.1 \\
rod tautness gate & 0.9 & 0.9 \\
pose time-out (s) & 1.0 & 0.25 \\
reference time-out (s) & 2.0 & 1.0 \\
planner solve budget (s) & none & 0.06 \\
\midrule
payload mass (kg) & \multicolumn{2}{c}{0.86} \\
height integral $k_i$, bound, gate & \multicolumn{2}{c}{0.4\,s$^{-1}$, 0.3\,m, 0.25\,m} \\
pretension, cap release (s) & \multicolumn{2}{c}{3.0, 2.0} \\
tracker initial-condition box & \multicolumn{2}{c}{0 (was 0.025 until 4 Oct 2026)} \\
\bottomrule
\end{tabularx}
\end{table}

\section{Development timeline}
\label{sec:app_timeline}

\begin{table}[H]
\centering
\small
\caption{Main steps of the work, 2026.}
\label{tab:app_timeline}
\begin{tabularx}{\textwidth}{lX}
\toprule
When & Step \\
\midrule
Feb--Apr & ROS~2 and Gazebo environment; four-UAV formation flight in simulation; one UAV with a tethered load \\
Jun & Cable models (spring--damper, then rigid rods); first port of the load--cable planner \\
Jul & Three UAVs carrying a payload with rigid rods in simulation; first rig flights; the dissipative controller and the first mid-flight detach \\
Aug & Experiment runner; OCP resize against network detach; attach integration and the weld-joint ladder; velocity tracking for the network \\
early Sep & Runtime attachment geometry and detach by OCP resize; network attach during a circle; layout studies \\
mid Sep & The 0.86\,kg ring in every model; measured rod lengths; thrust-gain trim; rod sweep from the floor; height integral; first carry attempts on the rig \\
late Sep & Detach layout 1/3/5/9; attach by OCP resize; integration with the collaborator's controllers (rejoin mission, ground join); rig days with the first lifts and the thrust identification; the rig twin \\
1--4 Oct & Three-UAV rig flights to 1.0\,m; the breakaway cap and its release; comparison of approaches to model error; the tracker initial-condition fix; integral action in the planner's model \\
\bottomrule
\end{tabularx}
\end{table}

\section{Negative results}
\label{sec:app_negatives}

Approaches that were built, tested and not kept. Each is recorded so that it is not repeated.

\begin{table}[H]
\centering
\small
\caption{Approaches measured and not kept.}
\label{tab:app_negatives}
\begin{tabularx}{\textwidth}{XX}
\toprule
Approach & Result \\
\midrule
% rows flown before the rig twin, listed in the box under the table:
% Terminal velocity reference in the tracker & slightly worse on all four tracking metrics \\
% Per-UAV integrators in the velocity loop & tilted the payload to 60$^\circ$ \\
% Network cone tilted onto the effective gravity & no radius benefit, formation contracted further \\
% Network tensions at the payload's measured attitude & weld transient 68$^\circ$, both runs aborted \\
% Horizontal leash on the payload's measured position & position loop closed; aborts, newcomer never welded \\
% Admittance on the measured rod tension & biased by the thrust model; worse hover \\
% Extra compliance at the weld (mid-span joint) & no change; a rigid weld capsized 5 of 5 \\
% Measured-force throttle (incremental inversion) & exact at the UAV, but stiff UAVs on approximate rods let rods go slack; removed \\
% Hand-back to an even ring after a detach & 62$^\circ$ tilt; survivors keep their real attachment points instead \\
% Load-force Kalman filter driven by the planned motion & converged to its own fixed point, not the target; added oscillation; removed \\
% Steeper nominal rods at the ground hand-over & same flip: the rigid rod turns any position mismatch into torque at the pivot \\
% Emulated rate-loop integral gain 10 & limit cycle at 3.1\,Hz; gain 5 adopted \\
Planner weight on the ring's horizontal position $\times 3$, $\times 10$ & push displacement cut only 23\,\% and 38\,\%; ring rocked at $\times 10$ \\
Rods aimed at the hand-over angle & push displacement worse (20.7 against 18.5\,cm) \\
Lumped force on the whole system from motion and thrust & blind to internal forces by Newton's law; on the rig it reads a steady external pull \\
\bottomrule
\end{tabularx}
\end{table}

\redo{the negative results flown before the rig twin, where they are still cited: a terminal velocity reference in the tracker; per-UAV integrators in the velocity loop; the network cone tilted onto the effective gravity; network tensions at the payload's measured attitude; a horizontal leash on the payload's measured position; an admittance on the measured rod tension; extra compliance at the weld; the measured-force throttle; the hand-back to an even ring after a detach; a load-force Kalman filter driven by the planned motion; steeper nominal rods at the ground hand-over; an emulated rate-loop integral gain of 10.}
